% ======================================================================
\section{Mesh-agnostic and operator-based representations}
\label{sec:meshfree}
% ======================================================================

Most methods discussed so far assume that every snapshot lives on the same
discretization, a structured grid in the case of CAEs. Engineering CFD relies on
unstructured, adapted, and moving meshes, and experimental data come on sensor
layouts that change between campaigns. Interpolating such data onto a Cartesian
grid degrades near-wall resolution and couples the learned representation to
an arbitrary choice of grid. This section reviews representations that relax
this assumption, and operator-learning architectures whose internal latent
spaces play the role of a reduced representation.

\subsection{Graph autoencoders}

Graph neural networks operate directly on mesh connectivity. The
encode--process--decode architecture of MeshGraphNets \parencite{pfaff2021meshgraphnets}
showed that message passing on simulation meshes can learn accurate dynamics
for a range of physical systems. It does not, however, reduce the number of
degrees of freedom. Graph autoencoders add pooling (graph coarsening) and
unpooling stages to obtain a genuine bottleneck. \textcite{gruber2022comparison}
found that a graph-convolutional AE is advantageous on data with irregular
connectivity when the latent space is sufficiently large.
\textcite{pichi2024gca} developed graph-convolutional AEs for model order
reduction of parametrized PDEs on unstructured meshes, and
\textcite{barwey2023multiscale} designed multiscale graph AEs whose adaptive
graph-reduction stages identify the spatial regions most relevant to the
dynamics, which provides a form of interpretability. Graph AEs are local by
construction, since each message-passing step aggregates information from a
neighbourhood. They are therefore the unstructured-mesh analogue of both CAEs
and the localized bases of \cref{sec:local}.

\subsection{Implicit neural representations}
\label{sec:inr}

An implicit neural representation (INR), or neural field, represents a flow
field as a coordinate network $\q(\bm{x})\approx g_{\omega}(\bm{x};\z)$
conditioned on a latent code $\z$. Because the decoder can be queried at any
point, the representation is independent of discretization. Encoding is done
either by an encoder network acting on point clouds or by ``auto-decoding'',
which optimizes $\z$ for each snapshot. Coordinate networks suffer severely
from spectral bias, which is mitigated by periodic activations
\parencite{sitzmann2020siren} or random Fourier features \parencite{tancik2020fourier}.

Neural implicit flow (NIF) \parencite{pan2023nif} introduced a mesh-agnostic
DR paradigm for spatio-temporal fluid data. A hypernetwork maps time and
parameters to the weights of a coordinate network, and the hypernetwork
bottleneck serves as the latent space. CROM \parencite{chen2023crom} uses an
INR decoder to build continuous reduced-order models whose latent dynamics
follow from the governing PDE, independent of the training discretization.
DINo \parencite{yin2023dino} evolves the INR latent code with a neural ODE,
giving forecasts that are continuous in space and time. CORAL
\parencite{serrano2023coral} performs operator learning entirely in INR latent
spaces on general geometries. AROMA \parencite{serrano2024aroma} replaces the
single global code with a set of spatially anchored latent tokens decoded by
\emph{local} neural fields and processed by a transformer, so the latent space
preserves spatial structure. This design echoes, in a learned setting, the
local-basis philosophy of DLS. For generative modeling of turbulence,
CoNFiLD \parencite{du2024cnf} trains a latent diffusion model on
conditional-neural-field latents, producing stochastic, spatiotemporally
coherent turbulent fields in irregular geometries.

\subsection{Neural operators and their implicit reduction}

Neural operators learn maps between function spaces
\parencite{kovachki2023neural} and are not usually described as DR methods, yet
most of them contain an explicit or implicit low-dimensional representation.
PCA-Net \parencite{bhattacharya2021pcanet} is explicit: it applies PCA to input
and output functions and learns a map between the coefficients. DeepONet
\parencite{lu2021deeponet} expresses the output as
$\sum_{j=1}^{p}b_j(\text{input})\,t_j(\bm{x})$. Its trunk network $t_j$ is a
learned set of basis functions and its branch network $b_j$ supplies their
coefficients, the same structure as a linear reduced basis with a nonlinear
coefficient map. The Fourier neural operator \parencite{li2021fno} truncates
each layer to a fixed number of Fourier modes, which is a linear spectral
reduction applied repeatedly inside a nonlinear network.
\textcite{kontolati2024latent} made the reduction explicit by training
DeepONets in the latent space of an autoencoder (L-DeepONet), which improved
accuracy and efficiency for high-dimensional, multiscale problems including
fluid flows. Several recent architectures learn \emph{adaptive}
reductions inside attention-based operators. The latent neural operator
\parencite{wang2024lno} maps inputs to a learned latent space by cross-attention.
Transolver \parencite{wu2024transolver} softly assigns mesh points to a small
number of learned ``physical-state slices'' and applies attention among them.
Universal physics transformers \parencite{alkin2024upt} encode arbitrary
particle or mesh data into a fixed set of latent tokens, propagate them, and
decode at arbitrary query points. The distinction between ``dimensionality
reduction followed by a dynamics model'' and ``end-to-end operator learning''
is therefore increasingly one of emphasis. The same questions of latent size,
locality, generalization, and accuracy saturation apply to both.
